\begin{pastexam}[label={px:ntu-112-bime-ds-06}]{臺大生機所 112 學年度（A）第 6 題}
（問答）

(16\%) What are the worst-time complexity of the \textbf{insert} and \textbf{delete-max}
operations, respectively, for \textbf{priority queue} of size N in the following four
implementations: (a) unordered linked list; (b) ordered array; (c) binary heap; and (d) binary
search tree.
\end{pastexam}

\begin{pastsoln}
答案：
\begin{center}
\begin{tabular}{lcc}
\toprule
實作 & insert & delete-max \\
\midrule
(a) unordered linked list & $O(1)$ & $O(N)$ \\
(b) ordered array & $O(N)$ & $O(1)$ \\
(c) binary heap & $O(\log N)$ & $O(\log N)$ \\
(d) binary search tree & $O(N)$ & $O(N)$ \\
\bottomrule
\end{tabular}
\end{center}

\begin{itemize}
\item \textbf{(a) 無序串列}：insert 把新節點接在串列前端，$O(1)$。 delete-max 不知道最大值在
哪裡，要走過全部 $N$ 個節點找出最大的，再刪掉（走的時候順便記住它的前一個節點，刪除 $O(1)$），
$\Theta(N)$。
\item \textbf{(b) 有序陣列}：由小到大排好，最大值在最後一格。 delete-max 直接拿走最後一格，
$O(1)$。 insert 要維持有序：找位置可以二分搜尋 $O(\log N)$，但插進去要把比它大的元素都往後挪
一格，最壞（新元素最小）挪 $N$ 個，$\Theta(N)$。
\item \textbf{(c) 二元堆積}（資結書第~8.1~節）：insert 放在最後再上浮，delete-max 把最後一個
元素搬到根再下沉，都只沿著一條根到葉的路徑走，完整二元樹的高度是 $\lfloor \lg N \rfloor$，
所以兩者都是 $\Theta(\log N)$。
\item \textbf{(d) 二元搜尋樹}（沒有保證平衡）：insert 沿一條路徑往下找位置；delete-max 從根一直
往右走到最右邊的節點再刪掉（它沒有右子節點，刪除很容易）。 兩者都是 $O(h)$，$h$ 是樹高。 最壞
情況下樹退化成一條鏈（例如依遞增順序插入），$h = N - 1$，兩者都是 $\Theta(N)$。 若改用平衡
搜尋樹（AVL 樹、紅黑樹，資結書第~10~章），$h = O(\log N)$，兩者都降為 $O(\log N)$。
\end{itemize}
\end{pastsoln}
